\documentclass[11pt,a4paper,twoside]{jarticle}
%==== 科研費LaTeX =============================================
%	2021(R03)年度 若手研究者海外挑戦プログラム（第１回）
%============================================================
% 2018-04-26: Taku: Made 2019 version.
%============================================================
\input{forms/form00_header}
% user01_header
%＝＝＝　様式のファイルの形式の指定　＝＝＝＝＝＝＝＝＝＝＝＝＝＝＝＝＝
%   PDFではなく、eps の様式を読み込む場合は、次の行の頭に「%」をつけてください。
\setboolean{usePDFform}{true}
%＝＝＝＝＝＝＝＝＝＝＝＝＝＝＝＝＝＝＝＝＝＝＝＝＝＝＝＝＝＝＝＝＝＝＝
\input{forms/form01_header}

% user02_header
%＝＝＝　予算の表の印刷　＝＝＝＝＝＝＝＝＝＝＝＝＝＝＝＝＝＝＝＝＝
% 予算の集計の表を出すためには、次の行の頭の%を消してください。
%\setboolean{BudgetSummary}{true}
%＝＝＝＝＝＝＝＝＝＝＝＝＝＝＝＝＝＝＝＝＝＝＝＝＝＝＝＝＝＝＝＝＝

%＝＝＝ For English, uncomment the next line to left-justify inside table columns.
%\renewcommand{\KLCLLang}{\KLCL}

% ＝＝＝　一部のページだけタイプセット　＝＝＝＝＝＝＝＝＝＝＝＝＝＝
% New in 2009 fall version!
% 選んだページだけタイプセットするには、次の例の頭の%を消し、並べてください。
% 複数のページを選ぶこともできます。
% 提出前には、必ず全てコメントアウト(頭に%をつける)してください。
%ーーーーーーーーーーーーーーーーーーーーーーーーーーーーーーーーー
%\KLTypesetPage{1}			% p.1 (or p.1を含む連続したページ),
%\KLTypesetPage{3}			% p.3 (or p.3を含む連続したページ),
%\KLTypesetPagesInRange{5}{6}	% p.5 ~ p.6,
%\KLTypesetPagesInRange{8}{10}	% and p.8 ~ p.10
%＝＝＝＝＝＝＝＝＝＝＝＝＝＝＝＝＝＝＝＝＝＝＝＝＝＝＝＝＝＝＝＝＝

% ===== my favorite packages ====================================
% ここに、自分の使いたいパッケージを宣言して下さい。
\usepackage{wrapfig}
% \usepackage{amssymb}
%\usepackage{mb}
% \usepackage{color} % でも科研費の書類はグレースケールで印刷されます
%\DeclareGraphicsRule{.tif}{png}{.png}{`convert #1 `dirname #1`/`basename #1 .tif`.png}
%==========================================================

\newcommand{\KLShouKeiLine}[1]{\cline{#1}}
%もし、小計の上の線を取れと事務に言われたら、
%「そのようなことは、記入要項に書かれていないし、学振はそのようなことは気にしていない。」と
% 突っぱねる。
% それでもなお消せと理不尽なことを言われたら、次の行の 最初の「%」を消す。	
%\renewcommand{\KLShouKeiLine}[1]{}

\newcommand{\KLBudgetTableFontSize}{small}	% 予算の表のフォントの大きさ: small, footnotesize
\newcommand{\KLFundsTableFontSize}{small}	%応募中、受入れ予定の研究費のフォントの大きさ：normalsize, small, footnotesize

% ===== my personal definitions ==================================
% ここに、自分のよく使う記号などを定義して下さい。
\newcommand{\klpionn}{K_L \to \pi^0 \nu \overline{\nu}}
\newcommand{\kppipnn}{K^+ \to \pi^+ \nu \overline{\nu}}


\input{forms/hook3} % for future maintenance
% ===== Global definitions for the PD form ======================
%　基本情報
%
%------ 研究課題名  -------------------------------------------
\newcommand{\研究課題名}{象の卵}

%----- 研究機関名と研究代表者の氏名-----------------------
\newcommand{\研究機関名}{逢坂大学}
\newcommand{\申請者氏名}{湯川秀樹}
\newcommand{\研究代表者氏名}{\申請者氏名}

%---- 研究期間の最終年度 ----------------
\newcommand{\研究期間の最終元号年度}{34}	%平成で、半角数字のみ
%=========================================================
\input{forms/form02_2021_header}	% <<<
\input{forms/form03_header}
\input{forms/form05_wakate_kaigai_1_header}
\input{forms/form07_header}
%============================================================
%endPrelude

\begin{document}
\input{forms/hook5} % for future maintenance
%============================================================
%     User Inputs
%============================================================

%form: wakate_kaigai_form_03.tex ; user: wakate_kaigai_03_past_and_publications.tex
%========== wakate_kaigai =========
%===== p. 03 現在までの研究状況 =============
\section{現在までの研究状況と研究成果}
\newcommand{\現在までの研究状況と研究成果}{%
%begin  現在までの研究状況と研究成果＝＝＝＝＝＝＝＝＝＝＝＝＝＝＝＝＝＝＝
	\noindent
	\textbf{(1) 現在までの研究状況}\\
	今までは、地球上で最大の生物、シロナガスクジラの卵の研究を進めようとしてきた。
	クジラの卵の場合は、高い水圧に耐える必要があるため、堅固の構造となっているはずであり、
	これが解明されれば、将来、深海潜水艇への応用も効く。
	しかし、シロナガスクジラの生息範囲が広い、海に潜っている時間が長い、
	生息数も減っている、などの原因により、
	卵を見つけることができなかった。
	
	そこで、\underline{地球で}最大の動物から、\underline{地上で}最大の動物に研究対象を変更する。
	\\
	
	\noindent
	\textbf{(2) 研究成果等}
	\begin{enumerate}
		\item \label{pub:theoegg} ``Theory of Elephant Eggs'', 
				\underline{Juzo Kara} {\it et al.},
				Phys.\ Rev.\ Lett. {\bf 800}, 800-804 (2005). 
				
		\item ``仔象は死んだ'', 
				\underline{Kobo Abe},
				安部公房全集, {\bf 26}, 100-200, (2004).
		\item ``The Elephant's Child (象の鼻はなぜ長い)'', 
				\underline{R.~Kipling},
				Nature, {\bf 999}, 777-779, (2003).
		\item ``You can't Lay an Egg If You're an Elephant'', 
				\underline{F.~Ehrlich},\\
				JofUR ({\tt www.universalrejection.org}), {\bf N/A}, N/A (2002).
		\item ``Egg of Elephant-Bird'', 
				\underline{A.~Cooper},
				Nature, {\bf 409}, 704-707 (2001).
				
	\end{enumerate}
%end  現在までの研究状況と研究成果 ＝＝＝＝＝＝＝＝＝＝＝＝＝＝＝＝＝＝＝＝
}

%form: wakate_kaigai_form_04.tex ; user: wakate_kaigai_04_purpose_and_reasons.tex
%========== wakate_kaigai =========
%===== p. 04 派遣先における研究計画等 =============
\section{派遣先における研究計画等}
\newcommand{\研究目的と内容と外国で研究することの意義}{%
%begin  研究目的と内容と外国で研究することの意義 (figureやtable使用可)＝＝＝＝＝＝＝＝＝＝＝＝＝＝＝＝＝＝＝
        \noindent
        \textbf{(1) 研究目的・内容}\\%
	象の卵の研究の目的は．．．
	\\
	
        \noindent
	\textbf{(2) 外国で研究することの意義}\\
	私は今まで、象の卵の可能性について主に文献を漁って研究をしてきた。
	そうした長年の研究の末分かったことの一つは、日本に現在、自然界に生息
	する象はいないということである。
	最も最近生息した象はケナガマンモスのようであるが、
	祖父が子供の頃には既に絶滅していたそうである。
	マンモスの氷漬けの個体は北海道で見つかったが、卵は見つかっていない。
	また最近では2005年に愛知県のある会場で氷漬けの個体が見つかったが、
	これは実は密かにロシアから持ち込まれたものであり、国産象ではない。
	
	こうした経験から、象の卵を日本で探していても見つからないということを
	強く実感し、海外で研究する決心をした次第である。
	特に、象の卵を探す夢を子供の頃に私に与えてくれた
	Dr. Seussにぜひとも指導を仰ぎたく、師の元に行って研究を行う。
%end  研究目的と内容と外国で研究することの意義 (figureやtable使用可) ＝＝＝＝＝＝＝＝＝＝＝＝＝＝＝＝＝＝＝＝
}

% hook9 : right before \end{document} ============

%endUserFiles
\input{forms/hook7} % for future maintenance

\input{forms/wakate_kaigai_1_forms}
%endFormatFile

\input{forms/hook9} % for future maintenance
\end{document}
